\documentclass[10pt,a4paper]{amsart}
\usepackage[margin=19mm]{geometry}
\usepackage{amsmath,amssymb,mathtools}
\usepackage[scr=rsfs,cal=euler]{mathalfa}
\usepackage{xfrac}
\usepackage[unicode,hidelinks,bookmarks=false]{hyperref}
\newcommand{\ie}{i.e.\ }
\newcommand{\cf}{cf.\ }
\newcommand{\resp}{respectively\ }
\newcommand{\Subsection}{\subsection}
\providecommand{\nobreakdash}{\nobreak\hbox{-}}
\setlength{\emergencystretch}{2em}
\multlinegap0pt
\usepackage{bm}
\usepackage{bbm}
\usepackage{dsfont}
\usepackage{xspace}
\usepackage{cancel}
\usepackage{mathtools}
\usepackage{amsthm}
\makeatletter
\@ifundefined{proposition}{\newtheorem{proposition}{Proposition}}{}
\makeatother
\title[Global weight optimization of frame structures under free-vi]{Global weight optimization of frame structures under free-vibration eigenvalue constraints}
\author{Marek Tyburec, Michal Ko\v{c}vara, Marouan Handa, Jan Zeman}
\date{Source: arXiv:2405.08894v3 (CC BY 4.0)}
\begin{document}
\maketitle

\noindent\emph{Source and licence.} Global weight optimization of frame structures under free-vibration eigenvalue constraints. Version arXiv:2405.08894v3 (CC BY 4.0), distributed under the Creative Commons Attribution 4.0 International licence, as recorded in the arXiv licence metadata for this exact version and shown on the abstract page https://arxiv.org/abs/2405.08894v3. Full rights notice: licenses/arxiv-2405.08894-CC-BY-4.0.txt.\par\smallskip

\noindent\textit{Editorial scope.} Counted unit: the Proposition whose source label is \texttt{prop:scaling} in the article (source0.tex lines 490--492, source label \texttt{prop:scaling}), delivered together with the article's own complete original proof of it (source0.tex lines 493--527, one \texttt{proof} environment of 35 lines). No mathematics is written, repaired or re-derived: statement and proof are the article's own text line for line. Required context retained uncounted: lemma:dynamic\_admissibility (lines 288--290); prop:exists\_gamma\_fv (lines 340--344); eq:RQs (lines 397--402); prop:monotonic (lines 408--460); eq:lamdiff (lines 410--459); eq:monotoneous (lines 450--456); prop:supremum (lines 465--482); eq:supremum (lines 467--469); eq:splitting (lines 470--481). Every definition, display and label of those lines is retained unchanged. Omitted from this package and not redistributed: 1--287, 291--339, 345--396, 403--407, 457--464, 482--489, 528--1527. Nothing inside a retained excerpt is omitted or altered: every retained body is delivered exactly as the article's lines are written, so no omission receipt of any kind accompanies this package. Citations: the retained excerpts carry no citation command at all, so no bibliography entry of the article is transcribed and no reference is redistributed. Reference adaptation: none was needed and none was made; every reference inside the delivered unit resolves inside it, so no unresolved reference or citation is produced. Printed slips kept literal: no printed word, symbol, hypothesis or number of the counted statement or of its proof was altered. Layout and class: the article's own class is \texttt{siamart220329} with packages including times, graphicx, amssymb, amsmath, mathtools, stanli, float, subcaption, pgfplots, color, bm, hyperref, cleveref, ulem; the wrapper is the lane's pinned \texttt{amsart} editorial wrapper at 10pt on A4, which restates the theorem-like environments the retained text needs and adds the article's own macro definitions that the retained text needs (none), with extra wrapper packages (bm, bbm, dsfont, xspace, cancel, mathtools). The delivered numbering is the unit's own. Assets: no retained span contains a figure environment, a graphics-inclusion command, a PDF-page inclusion, a table or a TikZ picture; no image file is copied into this package, the package declares no asset dependency and no image file is redistributed with it. Licence: version arXiv:2405.08894v3 of this article is distributed under the Creative Commons Attribution 4.0 International licence, as recorded in the arXiv licence metadata for this exact version and shown on the abstract page https://arxiv.org/abs/2405.08894v3. A full rights notice is delivered at licenses/arxiv-2405.08894-CC-BY-4.0.txt. Characters: every retained excerpt is pure ASCII, so the delivered engine, XeLaTeX with its default Latin Modern text font, reports no missing character.
\begin{proposition}[Mass admissibility]\label{lemma:dynamic_admissibility}
	Let $$\mathbf{\hat{a}}\in\left\{\mathbf{a} \; \vert \; \mathbf{K}_j(\mathbf{a}) - \underline{\lambda}_j \mathbf{M}_j(\mathbf{a}) \succeq 0, \mathbf{a}\ge \mathbf{0}\right\}.$$ Then, $\mathrm{Im}\left(\mathbf{M}_j(\hat{\mathbf{a}})\right) \subseteq \mathrm{Im}\left(\mathbf{K}_j(\hat{\mathbf{a}})\right)$.
\end{proposition}

\begin{proposition}[{Positive finite eigenvalue}]\label{prop:exists_gamma_fv}
	Let $\mathrm{Im}(\mathbf{M}_j(\mathbf{a})) \subseteq \mathrm{Im}(\mathbf{K}_j(\mathbf{a}))$. Then, there exists $\hat{\lambda}>0$ such that
	%
	$\mathbf{K}_j(\mathbf{a}) - \hat{\lambda} \mathbf{M}_j(\mathbf{a})\succeq 0$.
\end{proposition}

\begin{equation}\label{eq:RQs}
	\begin{multlined}
		\lambda_{\min,j}(\delta \tilde{\mathbf{a}}) =\\
		\inf_{\mathbf{v}_j \in \mathbb{R}^{n_{\mathrm{dof},j}}\setminus \mathrm{Ker}\left(\hat{\mathbf{M}}_j^{\langle 0\rangle}+\hat{\mathbf{M}}_j^{\langle 1\rangle}\right)} \frac{\mathbf{v}_j^\mathrm{T} \hat{\mathbf{K}}_j^{\langle 1\rangle}\mathbf{v}_j + \delta \mathbf{v}_j^\mathrm{T}\hat{\mathbf{K}}_j^{\langle 2\rangle} \mathbf{v}_j + \delta^2 \mathbf{v}_j^\mathrm{T}\hat{\mathbf{K}}_j^{\langle 3\rangle} \mathbf{v}_j}{\frac{1}{\delta}\mathbf{v}_j^\mathrm{T} \hat{\mathbf{M}}_j^{\langle 0\rangle} \mathbf{v}_j + \mathbf{v}_j^\mathrm{T} \hat{\mathbf{M}}_j^{\langle 1\rangle} \mathbf{v}_j}.
	\end{multlined}
\end{equation}

\begin{proposition}\label{prop:monotonic}
Let $\lambda_{\min,j}(\delta \tilde{\mathbf{a}})$ be as defined in \eqref{eq:RQs}. Then, $\lambda_{\min,j}(\delta \tilde{\mathbf{a}})$ is a non-decreasing function.
\begin{proof}
Let $\Delta > 0$; we need to show that $\lambda_{\min,j}((\delta+\Delta) \tilde{\mathbf{a}}) - \lambda_{\min,j}(\delta \tilde{\mathbf{a}})\ge 0$. Using the definition \eqref{eq:RQs}, we have %\textcolor{purple}{ We don't really need to write equations (3.11), it is clear enough to directly proceed to (3.12). }
\begin{multline}
	\lambda_{\min,j}((\delta + \Delta)\tilde{\mathbf{a}}) - \lambda_{\min,j}(\delta \tilde{\mathbf{a}})=\\
	%
	\inf_{\mathbf{v}_j \in \mathbb{R}^{n_{\mathrm{dof},j}}\setminus \mathrm{Ker}\left(\hat{\mathbf{M}}_j^{\langle 0\rangle}+\hat{\mathbf{M}}_j^{\langle 1\rangle}\right)} 
	\frac{\mathbf{v}_j^\mathrm{T} \hat{\mathbf{K}}_j^{\langle 1\rangle}\mathbf{v}_j + (\delta+\Delta) \mathbf{v}_j^\mathrm{T}\hat{\mathbf{K}}_j^{\langle 2\rangle} \mathbf{v}_j + (\delta+\Delta)^2 \mathbf{v}_j^\mathrm{T}\hat{\mathbf{K}}_j^{\langle 3\rangle} \mathbf{v}_j}{\frac{1}{\delta + \Delta }\mathbf{v}_j^\mathrm{T} \hat{\mathbf{M}}_j^{\langle 0\rangle} \mathbf{v}_j + \mathbf{v}_j^\mathrm{T} \hat{\mathbf{M}}_j^{\langle 1\rangle} \mathbf{v}_j} \\
	%
	- \inf_{\mathbf{v}_j \in \mathbb{R}^{n_{\mathrm{dof},j}}\setminus \mathrm{Ker}\left(\hat{\mathbf{M}}_j^{\langle 0\rangle}+\hat{\mathbf{M}}_j^{\langle 1\rangle}\right)} \frac{\mathbf{v}_j^\mathrm{T} \hat{\mathbf{K}}_j^{\langle 1\rangle}\mathbf{v}_j + \delta \mathbf{v}_j^\mathrm{T}\hat{\mathbf{K}}_j^{\langle 2\rangle} \mathbf{v}_j + \delta^2 \mathbf{v}_j^\mathrm{T}\hat{\mathbf{K}}_j^{\langle 3\rangle} \mathbf{v}_j}{\frac{1}{\delta}\mathbf{v}_j^\mathrm{T} \hat{\mathbf{M}}_j^{\langle 0\rangle} \mathbf{v}_j + \mathbf{v}_j^\mathrm{T} \hat{\mathbf{M}}_j^{\langle 1\rangle} \mathbf{v}_j}.
\end{multline}
%
Splitting the first infimum into two parts using the triangle inequality, we obtain 
%
\begin{multline}\label{eq:lamdiff}
	\lambda_{\min,j}((\delta + \Delta)\tilde{\mathbf{a}}) - \lambda_{\min,j}(\delta \tilde{\mathbf{a}})\ge\\
	%
	\inf_{\mathbf{v}_j \in \mathbb{R}^{n_{\mathrm{dof},j}}\setminus \mathrm{Ker}\left(\hat{\mathbf{M}}_j^{\langle 0\rangle}+\hat{\mathbf{M}}_j^{\langle 1\rangle}\right)} 
	\frac{\mathbf{v}_j^\mathrm{T} \hat{\mathbf{K}}_j^{\langle 1\rangle}\mathbf{v}_j + \delta \mathbf{v}_j^\mathrm{T}\hat{\mathbf{K}}_j^{\langle 2\rangle} \mathbf{v}_j + \delta^2 \mathbf{v}_j^\mathrm{T}\hat{\mathbf{K}}_j^{\langle 3\rangle} \mathbf{v}_j}{\frac{1}{\delta + \Delta}\mathbf{v}_j^\mathrm{T} \hat{\mathbf{M}}_j^{\langle 0\rangle} \mathbf{v}_j + \mathbf{v}_j^\mathrm{T} \hat{\mathbf{M}}_j^{\langle 1\rangle} \mathbf{v}_j} \\
	%
	+ \inf_{\mathbf{v}_j \in \mathbb{R}^{n_{\mathrm{dof},j}}\setminus \mathrm{Ker}\left(\hat{\mathbf{M}}_j^{\langle 0\rangle}+\hat{\mathbf{M}}_j^{\langle 1\rangle}\right)} 
	\frac{\Delta \mathbf{v}_j^\mathrm{T}\hat{\mathbf{K}}_j^{\langle 2\rangle} \mathbf{v}_j + (2\delta \Delta + \Delta^2) \mathbf{v}_j^\mathrm{T}\hat{\mathbf{K}}_j^{\langle 3\rangle} \mathbf{v}_j}{\frac{1}{\delta + \Delta}\mathbf{v}_j^\mathrm{T} \hat{\mathbf{M}}_j^{\langle 0\rangle} \mathbf{v}_j + \mathbf{v}_j^\mathrm{T} \hat{\mathbf{M}}_j^{\langle 1\rangle} \mathbf{v}_j}\\
	%
	- \inf_{\mathbf{v}_j \in \mathbb{R}^{n_{\mathrm{dof},j}}\setminus \mathrm{Ker}\left(\hat{\mathbf{M}}_j^{\langle 0\rangle}+\hat{\mathbf{M}}_j^{\langle 1\rangle}\right)} \frac{\mathbf{v}_j^\mathrm{T} \hat{\mathbf{K}}_j^{\langle 1\rangle}\mathbf{v}_j + \delta \mathbf{v}_j^\mathrm{T}\hat{\mathbf{K}}_j^{\langle 2\rangle} \mathbf{v}_j + \delta^2 \mathbf{v}_j^\mathrm{T}\hat{\mathbf{K}}_j^{\langle 3\rangle} \mathbf{v}_j}{\frac{1}{\delta}\mathbf{v}_j^\mathrm{T} \hat{\mathbf{M}}_j^{\langle 0\rangle} \mathbf{v}_j + \mathbf{v}_j^\mathrm{T} \hat{\mathbf{M}}_j^{\langle 1\rangle} \mathbf{v}_j}.
\end{multline}
%
Because $\forall \delta,\Delta >0: \frac{1}{\delta} \ge \frac{1}{\delta + \Delta}$, it holds for all $\mathbf{v}_j$, including the minimizing ones, that
\begin{equation}
	\frac{1}{\delta}\mathbf{v}_j^\mathrm{T} \hat{\mathbf{M}}_j^{\langle 0\rangle} \mathbf{v}_j \ge \frac{1}{\delta + \Delta}\mathbf{v}_j^\mathrm{T} \hat{\mathbf{M}}_j^{\langle 0\rangle} \mathbf{v}_j.
\end{equation}
%
Consequently,
% 
\begin{multline}
	\inf_{\mathbf{v}_j \in \mathbb{R}^{n_{\mathrm{dof},j}}\setminus \mathrm{Ker}\left(\hat{\mathbf{M}}_j^{\langle 0\rangle}+\hat{\mathbf{M}}_j^{\langle 1\rangle}\right)} 
	\frac{\mathbf{v}_j^\mathrm{T} \hat{\mathbf{K}}_j^{\langle 1\rangle}\mathbf{v}_j + \delta \mathbf{v}_j^\mathrm{T}\hat{\mathbf{K}}_j^{\langle 2\rangle} \mathbf{v}_j + \delta^2 \mathbf{v}_j^\mathrm{T}\hat{\mathbf{K}}_j^{\langle 3\rangle} \mathbf{v}_j}{\frac{1}{\delta + \Delta}\mathbf{v}_j^\mathrm{T} \hat{\mathbf{M}}_j^{\langle 0\rangle} \mathbf{v}_j + \mathbf{v}_j^\mathrm{T} \hat{\mathbf{M}}_j^{\langle 1\rangle} \mathbf{v}_j} \\
	- \inf_{\mathbf{v}_j \in \mathbb{R}^{n_{\mathrm{dof},j}}\setminus \mathrm{Ker}\left(\hat{\mathbf{M}}_j^{\langle 0\rangle}+\hat{\mathbf{M}}_j^{\langle 1\rangle}\right)} \frac{\mathbf{v}_j^\mathrm{T} \hat{\mathbf{K}}_j^{\langle 1\rangle}\mathbf{v}_j + \delta \mathbf{v}_j^\mathrm{T}\hat{\mathbf{K}}_j^{\langle 2\rangle} \mathbf{v}_j + \delta^2 \mathbf{v}_j^\mathrm{T}\hat{\mathbf{K}}_j^{\langle 3\rangle} \mathbf{v}_j}{\frac{1}{\delta}\mathbf{v}_j^\mathrm{T} \hat{\mathbf{M}}_j^{\langle 0\rangle} \mathbf{v}_j + \mathbf{v}_j^\mathrm{T} \hat{\mathbf{M}}_j^{\langle 1\rangle} \mathbf{v}_j} \ge 0,
\end{multline}
%
which allows us to simplify \eqref{eq:lamdiff} to
%
\begin{equation}\label{eq:monotoneous}
\begin{multlined}
\lambda_{\min,j}((\delta + \Delta)\tilde{\mathbf{a}}) - \lambda_{\min,j}(\delta \tilde{\mathbf{a}}) \ge\\
 \inf_{\mathbf{v}_j \in \mathbb{R}^{n_{\mathrm{dof},j}}\setminus \mathrm{Ker}\left(\hat{\mathbf{M}}_j^{\langle 0\rangle}+\hat{\mathbf{M}}_j^{\langle 1\rangle}\right)} 
\frac{\Delta \mathbf{v}_j^\mathrm{T}\hat{\mathbf{K}}_j^{\langle 2\rangle} \mathbf{v}_j + (2\delta \Delta + \Delta^2) \mathbf{v}_j^\mathrm{T}\hat{\mathbf{K}}_j^{\langle 3\rangle} \mathbf{v}_j}{\frac{1}{\delta + \Delta}\mathbf{v}_j^\mathrm{T} \hat{\mathbf{M}}_j^{\langle 0\rangle} \mathbf{v}_j + \mathbf{v}_j^\mathrm{T} \hat{\mathbf{M}}_j^{\langle 1\rangle} \mathbf{v}_j} \ge 0.
\end{multlined}
\end{equation}
%
In \eqref{eq:monotoneous}, the nonnegativity follows from $\Delta \mathbf{v}_j^\mathrm{T}\hat{\mathbf{K}}_j^{\langle 2\rangle} \mathbf{v}_j \ge 0$ and $\left(2\delta \Delta + \Delta^2\right) \mathbf{v}_j^\mathrm{T}\hat{\mathbf{K}}_j^{\langle 3\rangle} \mathbf{v}_j\ge 0$, which is because of $\hat{\mathbf{K}}_j^{\langle 2\rangle}, \hat{\mathbf{K}}_j^{\langle 3\rangle} \succeq 0$ and $\delta\ge 0, \Delta > 0$.
\end{proof}
\end{proposition}

\begin{proposition}\label{prop:supremum}
Let $\lambda_{\min,j}(\delta \tilde{\mathbf{a}})$ be as defined in \eqref{eq:RQs}, and let $\mathbf{v}_j := \mathbf{t}_j + \mathbf{w}_j$ as defined above. Then, 
\begin{equation}\label{eq:supremum}
	\sup_{\delta} \lambda_{\min,j}(\delta \tilde{\mathbf{a}}) = \inf_{\mathbf{w}_j \in \mathrm{Ker}\left(\hat{\mathbf{K}}_j^{\langle 2\rangle} + \hat{\mathbf{K}}_j^{\langle 3\rangle}\right) \cap \mathrm{Im}\left(\hat{\mathbf{M}}_j^{\langle 0\rangle} + \hat{\mathbf{M}}_j^{\langle 1\rangle}\right)} \frac{\mathbf{w}_j^\mathrm{T} \hat{\mathbf{K}}_j^{\langle 1\rangle}\mathbf{w}_j}{\mathbf{w}_j^\mathrm{T} \hat{\mathbf{M}}_j^{\langle 1\rangle} \mathbf{w}_j}.
\end{equation}
\begin{proof}
We prove the statement by splitting $\mathbf{v}_j$ into $\mathbf{t}_j$ and $\mathbf{w}_j$ and showing that it must hold that $\mathbf{t}_j=\mathbf{0}$.  For this splitting, the supremum is expressed as
\begin{equation}\label{eq:splitting}
	\begin{multlined}
		\sup_\delta \lambda_{\min,j}(\delta \tilde{\mathbf{a}})= \\
		\sup_\delta \inf\limits_{\mathbf{t}_j, \mathbf{w}_j}
		 \frac{\left(\mathbf{t}_j+\mathbf{w}_j\right)^\mathrm{T} \hat{\mathbf{K}}_j^{\langle 1\rangle}\left(\mathbf{t}_j+\mathbf{w}_j\right) + \delta\mathbf{t}_j^\mathrm{T} \hat{\mathbf{K}}_j^{\langle 2\rangle}\mathbf{t}_j + \delta^2\mathbf{t}_j^\mathrm{T} \hat{\mathbf{K}}_j^{\langle 3\rangle}\mathbf{t}_j }{\frac{1}{\delta}\left(\mathbf{t}_j+\mathbf{w}_j\right)^\mathrm{T} \hat{\mathbf{M}}_j^{\langle 0\rangle} \left(\mathbf{t}_j+\mathbf{w}_j\right) + \left(\mathbf{t}_j+\mathbf{w}_j\right)^\mathrm{T} \hat{\mathbf{M}}_j^{\langle 1\rangle} \left(\mathbf{t}_j+\mathbf{w}_j\right)}.
	\end{multlined}
\end{equation}
%
Because of the monotonicity (recall Proposition~\ref{prop:monotonic}), $\sup_\delta \lambda_{\min,j}(\delta \tilde{\mathbf{a}}) = \lim\limits_{\delta\rightarrow \infty}\lambda_{\min,j}(\delta \tilde{\mathbf{a}})$. Further, suppose by {contradiction} that $\mathbf{t}_j \neq \mathbf{0}$ holds at the infimum. Then, $\lambda_{\min,j}(\delta \tilde{\mathbf{a}}) \rightarrow \infty$ as $\delta \rightarrow \infty$, due to $\mathbf{w}_j \in \mathrm{Im}\left(\hat{\mathbf{M}}_j^{\langle 0\rangle} + \hat{\mathbf{M}}_j^{\langle 1\rangle}\right)$ and because $\mathbf{t}_j$ cannot be in $\mathrm{Ker}\left(\hat{\mathbf{K}}_j^{\langle 2\rangle} + \hat{\mathbf{K}}_j^{\langle 3\rangle}\right)$. On the other hand, for $\mathbf{t}_j = \mathbf{0}$, $\lambda_{\min,j}(\delta \tilde{\mathbf{a}})$ is finite for all $\delta$ due to $\mathbf{w}_j \in \mathrm{Im}\left(\hat{\mathbf{M}}_j^{\langle 0\rangle} + \hat{\mathbf{M}}_j^{\langle 1\rangle}\right)$. Consequently, $\mathbf{t}_j^*=\mathbf{0}$ must hold at the infimum. After inserting $\mathbf{t}_j^*=\mathbf{0}$ and $\delta^*\rightarrow \infty$ into \eqref{eq:splitting}, we receive \eqref{eq:supremum}, which completes the proof.
\end{proof}
\end{proposition}

\begin{proposition}\label{prop:scaling}
	There exists $\delta$ such that $\mathbf{K}_j(\delta \tilde{\mathbf{a}}) - \underline{\lambda}_j \mathbf{M}_j(\delta \tilde{\mathbf{a}}) \succeq 0$ if and only if $\tilde{\mathbf{a}} \in \left\{\mathbf{a} \;\vert\; \mathbf{a}\ge \mathbf{0}, \mathrm{Im}\left(\mathbf{M}_j(\mathbf{a})\right) \subseteq \mathrm{Im}\left(\mathbf{K}_j(\mathbf{a})\right),\mathbf{W}_j^\mathrm{T}\left[\hat{\mathbf{K}}_{j}^{\langle 1\rangle} - \underline{\lambda}_j \hat{\mathbf{M}}_{j}^{\langle 1\rangle} \right]\mathbf{W}_j \succeq 0\right\}$, where $\mathbf{W}_j$ denotes an orthonormal basis of $\mathrm{Ker}\left(\hat{\mathbf{K}}_j^{\langle 2\rangle} + \hat{\mathbf{K}}_j^{\langle 3\rangle}\right) \cap \mathrm{Im}\left(\hat{\mathbf{M}}_j^{\langle 0\rangle} + \hat{\mathbf{M}}_j^{\langle 1\rangle}\right)$.
\end{proposition}

\begin{proof}
\fbox{$\Rightarrow$} Let $\mathbf{K}_j(\delta \tilde{\mathbf{a}}) - \underline{\lambda}_j \mathbf{M}_j(\delta \tilde{\mathbf{a}}) \succeq 0$ hold. Then, $\mathrm{Im}\left(\mathbf{M}_j(\tilde{\mathbf{a}})\right) \subseteq \mathrm{Im}\left(\mathbf{K}_j(\tilde{\mathbf{a}})\right)$ based on Proposition \ref{lemma:dynamic_admissibility}. Next, we project the matrix inequality using $\mathbf{W}_j$, yielding
%
\begin{equation}
	\delta \mathbf{W}_{{j}}^\mathrm{T} \hat{\mathbf{K}}_j^{\langle 1\rangle} \mathbf{W}_j- \underline{\lambda}_j \delta \mathbf{W}_{{j}}^\mathrm{T} \hat{\mathbf{M}}_j^{\langle 1\rangle} \mathbf{W}_j - \underline{\lambda}_j\mathbf{W}_{{j}}^\mathrm{T} \hat{\mathbf{M}}_j^{\langle 0\rangle} \mathbf{W}_j \succeq 0.
\end{equation}
%
For $\delta>0$, this is further equivalent to
%
\begin{equation}
	\mathbf{W}_{{j}}^\mathrm{T} \hat{\mathbf{K}}_j^{\langle 1\rangle} \mathbf{W}_j- \underline{\lambda}_j \mathbf{W}_{{j}}^\mathrm{T} \hat{\mathbf{M}}_j^{\langle 1\rangle} \mathbf{W}_j \succeq \frac{1}{\delta}\underline{\lambda}_j\mathbf{W}_{{j}}^\mathrm{T} \hat{\mathbf{M}}_j^{\langle 0\rangle} \mathbf{W}_j \succeq 0.
\end{equation}
%
\fbox{$\Leftarrow$} Let $\mathrm{Im}\left(\mathbf{M}_j(\tilde{\mathbf{a}})\right) \subseteq \mathrm{Im}\left(\mathbf{K}_j(\tilde{\mathbf{a}})\right)$ and $\mathbf{W}_j^\mathrm{T}\left[\hat{\mathbf{K}}_{j}^{\langle 1\rangle} - \underline{\lambda}_j \hat{\mathbf{M}}_{j}^{\langle 1\rangle} \right]\mathbf{W}_j \succeq 0$ hold. Because of the first assumption and Proposition \ref{prop:exists_gamma_fv}, it holds that
%
\begin{equation}
\forall \mathbf{v}_j \in \mathrm{Im}\left(\mathbf{K}_j(\tilde{\mathbf{a}})\right) \exists \gamma>0: \mathbf{v}_j^\mathrm{T}\mathbf{K}_j(\tilde{\mathbf{a}})\mathbf{v}_j - \gamma\mathbf{v}_j^\mathrm{T}\mathbf{M}_j(\tilde{\mathbf{a}})\mathbf{v}_j \ge 0,
\end{equation}
%
and thus also $\mathbf{K}_j(\tilde{\mathbf{a}}) - \gamma\mathbf{M}_j(\tilde{\mathbf{a}}) \succeq 0$ for $\gamma > 0$. Consequently, it remains to show that $\gamma\ge \underline{\lambda}_j$.

Because of $\mathbf{W}_j^\mathrm{T}\left[\hat{\mathbf{K}}_{j}^{\langle 1\rangle} - \underline{\lambda}_j \hat{\mathbf{M}}_{j}^{\langle 1\rangle} \right]\mathbf{W}_j \succeq 0$, we also have
%
\begin{equation}\label{eq:inequality}
	\forall \mathbf{w}_j \in \mathrm{Ker}\left(\hat{\mathbf{K}}_j^{\langle 2\rangle} + \hat{\mathbf{K}}_j^{\langle 3\rangle}\right) \cap \mathrm{Im}\left(\hat{\mathbf{M}}_j^{\langle 0\rangle} + \hat{\mathbf{M}}_j^{\langle 1\rangle}\right):  \mathbf{w}_j^\mathrm{T} \left[\hat{\mathbf{K}}_{j}^{\langle 1\rangle} - \underline{\lambda}_j \hat{\mathbf{M}}_{j}^{\langle 1\rangle} \right] \mathbf{w}_j \ge 0.
\end{equation}
%
In particular, the above holds for $w_j$ minimizing the left-hand side of \eqref{eq:inequality}, yielding
%
\begin{equation}
	\underline{\lambda}_j \le \inf_{\mathbf{w}_j \in \mathrm{Ker}\left(\hat{\mathbf{K}}_j^{\langle 2\rangle} + \hat{\mathbf{K}}_j^{\langle 3\rangle}\right) \cap \mathrm{Im}\left(\hat{\mathbf{M}}_j^{\langle 0\rangle} + \hat{\mathbf{M}}_j^{\langle 1\rangle}\right)} \frac{\mathbf{w}_j^\mathrm{T} \hat{\mathbf{K}}_{j}^{\langle 1\rangle} \mathbf{w}_j}{\mathbf{w}_j^\mathrm{T} \hat{\mathbf{M}}_{j}^{\langle 1\rangle} \mathbf{w}_j} = \sup_{\delta} \lambda_{\min,j}(\mathbf{a}(\delta))
\end{equation}
%
based on Proposition \ref{prop:supremum}.
\end{proof}

\end{document}
